%! TeX root = ../proposal_magang.tex
\chapter{TIMELINE KEGIATAN}
\label{chap:proposal-timeline-kegiatan}

\paragraphtextstyle

Untuk memberi gambaran kepada perusahaan / instansi tempat magang mengenai kegiatan-kegiatan apa saja yang akan dilaksanakan selama magang, maka peserta magang perlu menuliskan rencana pelaksanaan kegiatan yang diusulkan. Umumnya rencana pelaksanaan kegiatan dimulai dengan perkenalan tempat magang yang kemudian dilanjutkan sesuai dengan bidang / topik usulan magang.

\begingroup
\fontsize{10}{12}\selectfont
\setlength{\tabcolsep}{1.5pt}
\renewcommand{\arraystretch}{1.15}
\begin{longtable}{|>{\centering\arraybackslash}p{7mm}|>{\raggedright\arraybackslash}p{53mm}|*{9}{>{\centering\arraybackslash}p{7mm}|}}
\caption{\textbf{\textit{Rekap Rencana Kegiatan Magang [Periode Pelaksanaan]}}}
\label{tab:proposal-kegiatan-magang}\\
\hline
\textbf{No.} & \centering\arraybackslash\textbf{Uraian Kegiatan Magang} & \textbf{P1} & \textbf{P2} & \textbf{P3} & \textbf{P4} & \textbf{P5} & \textbf{P6} & \textbf{P7} & \textbf{P8} & \textbf{P9} \\
\hline
\endfirsthead
\hline
\textbf{No.} & \centering\arraybackslash\textbf{Uraian Kegiatan Magang} & \textbf{P1} & \textbf{P2} & \textbf{P3} & \textbf{P4} & \textbf{P5} & \textbf{P6} & \textbf{P7} & \textbf{P8} & \textbf{P9} \\
\hline
\endhead
1 & [Deskripsikan kegiatan utama, langkah pelaksanaan, dan target.] & \cellcolor[HTML]{B6D7A8} & & & & & & & & \\
\hline
2 & [Tuliskan kegiatan berikutnya beserta metode yang digunakan.] & & \cellcolor[HTML]{B6D7A8} & \cellcolor[HTML]{B6D7A8} & & & & & & \\
\hline
3 & [Jelaskan proses pelaksanaan dan hasil yang diharapkan.] & \cellcolor[HTML]{B6D7A8} & \cellcolor[HTML]{B6D7A8} & & & & & & & \\
\hline
\end{longtable}
\endgroup
